\section{Introduction}

Finite state automata are mathematical models of computation.
They are mostly used to study regular languages, i.e. languages with a certain structure.
This structure is also captured by regular expressions.

In this report we discuss our Haskell-implementation of finite state automata and regular expressions.
In particular, we implemented deterministic finite state automata, non-deterministic finite state automata and regular expressions.
The goal is to make the equivalence of finite state automata and regular languages explicit by various translation functions. 
Regular expressions are suitable for generating words from a regular language, but finite automata are more suitable for determining whether 
a word is in its language. Our implementation enables combining these features. For example, given a regular expression, we can compute another one that generates
the complement language using finite state automata. Moreover, given an automaton, we can generate words that it accepts using regular expressions. Furthermore,
we provide a toolkit that determines the minimal equivalent DFA to a DFA and checks isomorphism of DFAs.

Our report is structured as follows. In Section \ref{theory} we give a brief introduction to the theory of regular languages and point towards more detailed references. 
In Sections \ref{sec:DFA} - \ref{sec:RegExp} we discuss our implementation of finite state automata and regular expressions. In Section \ref{sec:Trans} we turn to the
implementation of the translation functions. We then provide extensive testing apparatus for our implementation in Section \ref{sec:simpletests} and conclude in Section \ref{sec:Conclusion}. 